\documentclass[11pt, reqno]{article}

\usepackage{amsmath, amsthm, amssymb}
\usepackage{enumitem}
\usepackage{tcolorbox}
\usepackage{hyperref}
\usepackage{tikz}
\usepackage{tikz-cd}
\usepackage{pgfplots}
\pgfplotsset{compat=1.18}
\usetikzlibrary{arrows.meta}
\usepackage{mathrsfs}
\usepackage{fancyhdr}
\usepackage[bottom=0.75in, top=1in, left=0.5in, right=0.5in]{geometry}
\usepackage{array}   % for \newcolumntype macro
\newcolumntype{L}{>{$}l<{$}}

\theoremstyle{plain}
\newtheorem*{theorem}{Theorem}
\newtheorem*{proposition}{Proposition}
\newtheorem{exercise}{Exercise}
\newtheorem*{lemma}{Lemma}
\newtheorem*{corollary}{Corollary}

\theoremstyle{definition}
\newtheorem*{definition}{Definition}
\newtheorem*{example}{Example}

\theoremstyle{remark}
\newtheorem*{remark}{Remark}

\renewcommand{\phi}{\varphi}
\renewcommand{\epsilon}{\varepsilon}
\renewcommand{\emptyset}{\varnothing}

\newcommand{\RR}{\mathbb{R}}
\newcommand{\ZZ}{\mathbb{Z}}
\newcommand{\NN}{\mathbb{N}}
\newcommand{\CC}{\mathbb{C}}
\newcommand{\QQ}{\mathbb{Q}}

\DeclareMathOperator{\ima}{\text{im}}

\begin{document}

\topmargin=-40pt
\rhead{Henry Woodburn}
\lhead{Math 622}
\renewcommand{\headrulewidth}{1pt}
\renewcommand{\headsep}{20pt}
\thispagestyle{fancy}

{\Huge \bfseries \noindent Homework 10}

\begin{enumerate}
    \item[1.] Let $U$ be the positive octant of $\RR^3$. Let $D$ be the distribution spanned 
    by the vector fields 
    \[
        X = y\frac{\partial}{\partial z} - z \frac{\partial}{\partial y}, \quad
        Y = z \frac{\partial}{\partial x} - x \frac{\partial}{\partial z}.
    \]
    (a.) We need to show $D$ is involutive. It suffices to verify that the Lie bracket
    $[X, Y]$ is contained in $D$, since any element of $D$ is a linear combination of $X$ and $Y$, 
    and by the bilinearity of the Lie bracket. We calculate that
    \[
        [X, Y] = y \frac{\partial}{\partial x} - x \frac{\partial}{\partial y} = \frac{x}{z}X + \frac{y}{z}Y.
    \]
    These coefficients are smooth functions on the open set $U$.

    (b.)  Notice both $X$ and $Y$ are perpendicular to the vector field $A = x \frac{\partial}{\partial x}
    + y \frac{\partial}{\partial y} + z \frac{\partial}{\partial z}$, since for example, at $p \in U$,
    \[
        X_p \cdot A_p = 0.
    \]  
    Then So $D_p$ is orthogonal to the vector $p$. But $D_p$ is two dimensional, so it is the entire orthogonal
    compliment of $p$. Then any integral submanifold of $D$ must have its tangent bundle orthogonal to the 
    radial vector field $A$. This is true for a sphere of any radius, restricted to $U$. 
    
    Moreover, any submanifold which has does not lie in a sphere must have varying radius. Suppose
    $M$ is such a submanifold of $U$, and choose points $p$ and $q$ with different distances from the origin.
    Consider a path from $p$ to $q$ lying in the plane defined by the points $p$, $q$, and the origin. 
    After converting to ordinary $2$-d polar coordinates on this plane, by the mean value theorem, 
    the velocity vector of this path must have non-zero radial component for some $t$, say at a point $p'$. 
    But this velocity vector must lie in the tangent space of $M$, and thus $M$ is not an integral manifold.
    
    Then the only integral submanifolds of $D$ are subsets of spheres. A maximal integral submanifold
    is thus an entire octant of a sphere of nonzero radius. 

    \item[2.] Let $D$ be the distribution of $\RR^3$ spanned by 
    \[
        X = \frac{\partial}{\partial x} + yz \frac{\partial}{\partial z} \quad Y = \frac{\partial}{\partial y}.
    \]
    (a.) An integral submanifold of $D$ passing through the origin is the plane $z = 0$. The 
    $z$ component of $X$ vanishes when $z = 0$, and its tangent space is clearly spanned by
    $\frac{\partial}{\partial x}$ and $\frac{\partial}{\partial y}$. 

    (b.) We can compute the bracket
    \[
        [X, Y] = z\frac{\partial}{\partial z},
    \]
    which does not lie in the span of $X$ and $Y$, as there is no way to eliminate either the $x$ or 
    the $y$ components by taking linear combinations. This demonstrates that a distribution may have integral
    submanifolds at some locations, but still not be involutive.

    \item[3.] Let $I$ be an ideal of forms on $M$ locally generated by $\{\omega^i\}_1^r$ on an open 
    set $U$.

    First suppose $I$ is a differential ideal. Let $\omega := \omega^1 \wedge \cdots \wedge \omega^r$. 
    Using the product rule for forms, we can iterate
    \begin{align*}
        d\omega = d(\omega^1 \wedge \cdots \wedge \omega^r) = 
        d\omega^1 \wedge (\omega^2 \wedge \cdots \wedge \omega^r) + 
        \omega^1 \wedge d(\omega^2 \wedge \cdots \wedge \omega^r) = \\
        \vdots\\
        d\omega^1 \wedge \omega^2 \wedge \cdots \wedge \omega^r + 
        \omega^1 \wedge d\omega^2 \wedge \cdots \wedge \omega^r + 
        \cdots + \omega^1 \wedge \cdots \wedge d\omega^r.
    \end{align*}
    Then using the fact that $I$ is a differential ideal, write $d\omega^i = \sum \alpha_j^i \wedge \omega^j$
    for each $i$. Then many of the terms above vanish, leaving only
    \[
        \sum \alpha_i^i \wedge \omega = \alpha \wedge \omega
    \]
    where $\alpha = \sum \alpha_i^i$.

    Conversely suppose $d\omega = \alpha \wedge \omega$ for some $\alpha$. We can complete the forms to a basis $\{\omega^k\}_1^n$ by restricting $U$. For each $i$ write
    $d\omega^i = \sum_{I} a_I^i \omega^I$, where $I$ ranges over all $2$ element subsets of $[1,\dots, n]$. Then
    \begin{equation}\label{1}
        d\omega = \sum_i (-1)^{i-1} d\omega^i \wedge (\omega^1 \wedge \cdots \wedge \widehat{\omega^i} \wedge \cdots \wedge \omega^r)
        = \sum_{I, i} a_I^i \omega^I \wedge (\omega^1 \wedge \cdots \wedge \widehat{\omega^i} \wedge \cdots \wedge \omega^r).
    \end{equation}
    
    
    Also write $\alpha = \sum_1^n \alpha_i \omega^i$. Then we have
    \begin{equation}\label{2}
        d\omega = \sum_1^n \alpha_i \omega^i \wedge \omega = \sum_{r + 1}^n \alpha_i \omega^i \wedge \omega.
    \end{equation}
    Then the sums \ref{1} and \ref{2} are equal. Both are composed of basis elements of the $r + 1$ forms induced by 
    the $\omega^i$'s, so by linear independence, each $a_I^i$ for which both elements of $I$ are greater than
    $r$ must be zero. Then every term in the expressions for the $d\omega^i$ contains an element of the 
    generating set $\{\omega^1,\dots, \omega^r\}$, and can be written in the form
    \[
        d\omega^i = \sum_1^r \beta_i \wedge \omega^i
    \]
    for some $\beta_i$. This proves $I$ is differential for the generating forms $\omega^i$. For 
    arbitrary elements of $I$, we can prove their differential is in $I$ by writing them as combinations
    of wedges with basis elements and expanding using the definition of the differential.

    \item[4.] Let $D$ be a smooth distribution of rank $k$ on $M$ an $n$ dimensional manifold. Let $\omega^1,
    \dots, \omega^{n-k}$ be its defining forms within a set $U$. 

    First suppose $D$ is involutive. The set of forms which vanish on $D$ is an ideal $I_D$. It is generated
    by the defining forms $\omega^1, \dots, \omega^{n-k}$. From the notes, we know that $D$ is involutive
    if and only if $I_D$ is differential. Then we can write 
    \[
        d\omega^i = \sum_{j = 1}^{n-k} \alpha_j^i \wedge \omega j
    \]
    for some $\alpha_j^i$. Then it follows that 
    \[
        d\omega^i\wedge\omega^1 \wedge \cdots \wedge \omega^{n-k} = 0
    \]
    after expanding $d\omega^i$, since all terms contain duplicate indices. 

    Conversely suppose $d\omega^i \wedge \omega^1 \wedge \cdots \wedge \omega^{n-k} = 0$ for all $i = 1, \dots, n-k$.
    It suffices to prove $d\omega^i \in I_D$ for all $i$ by a theorem in Lee. If we expand $d\omega^i$
    after completing our set to a basis, there cannot be any terms which do not contain some $\omega^i$
    for $i \in [1, \dots, n-k]$,
    otherwise $d\omega^i \wedge \omega^1 \wedge \cdots \wedge \omega^{n-k}$ would not vanish. 

    \item[6.] Let $D$ be a distribution on $M$ and $X$ a vector field, and $e^{tX}$ the local flow of $X$. 
    (a.) suppose $(e^{tX})_* D = D$ for small $t$. The Lie derivative is defined as 
    \[
        \mathcal{L}_X(Y) = \frac{d}{dt}(e^{tX})_* Y|_{t = 0} = [X, Y].
    \]
    The condition above guarantees that the time derivative will also be in $D$, by writing as a difference quotient.
    So $[X,Y] \in D$ for all $Y \in D$. 

    (b.) Suppose $[X, Y] \in D$ when $Y \in D$. If $e^{tX}$ is defined, for $Y \in D$, we have 
    \[
        [X,Y] = \frac{d}{dt}(e^{tX})_* Y|_{t = 0}.
    \]
    Let $Y_t = (e^{tX})_* Y$. However by the property of local flows we can take the restriction at any time close
    to zero. Then the problem is to find a flow for which $[X,Y]$ is the infinitesimal generator. But 
    this can be solved locally within $D$ since $[X,Y] \in D$. Thus $D$ is invariant under pullback of $Y$ 
    by the flow of $X$.

    (c.) That $[X,Y]$ is an infinitesimal symmetry of $D$ when $X$ and $Y$ are follows from writing 
    \[
        [Z,[X,Y]] = -[X, [Y,Z]] - [Y, [Z,X]],
    \]
    where $Z \in D$, and noting first that the inner brackets are in $D$, then the outer ones.


\end{enumerate}

\end{document}